% !TeX program = lualatex
% Compile twice: lualatex open_problems_500.tex
% All references and prose are embedded; no BibTeX database or external inputs.
\documentclass[11pt,a4paper]{article}
\usepackage[margin=24mm,headheight=14pt]{geometry}
\usepackage{fontspec}
\setmainfont{Latin Modern Roman}
\setsansfont{Latin Modern Sans}
\setmonofont{Latin Modern Mono}
\usepackage{amsmath,amssymb,mathtools}
\usepackage{unicode-math}
\setmathfont{Latin Modern Math}
\usepackage{microtype}
\usepackage{enumitem,xurl,needspace}
\usepackage[dvipsnames]{xcolor}
\usepackage{hyperref}
\hypersetup{unicode=true,colorlinks=true,linkcolor=MidnightBlue,urlcolor=MidnightBlue,
 pdftitle={Research notebook: Problems 40--49},pdfauthor={Source-annotated compilation},bookmarksnumbered=true}
\usepackage{bookmark}
\usepackage{tocloft}
\setlength{\cftsecnumwidth}{3em}
\cftsetpnumwidth{2.5em}
\cftsetrmarg{4em}
\usepackage{titlesec}
\titleformat{\section}{\Large\bfseries\raggedright}{\thesection}{0.7em}{}
\usepackage{fancyhdr}
\pagestyle{fancy}\fancyhf{}
\fancyhead[L]{\small\sffamily Mathematical problem statements}
\fancyhead[R]{\small Source edition 22}
\fancyfoot[C]{\thepage}
\setlength{\parindent}{0pt}
\setlength{\parskip}{5pt plus 1pt}
\setlength{\emergencystretch}{3em}
\setcounter{tocdepth}{1}
\setcounter{secnumdepth}{1}
\setlist[itemize]{leftmargin=3em,itemsep=5pt,topsep=4pt,parsep=0pt}
\allowdisplaybreaks
\newcommand{\N}{\mathbb{N}}
\newcommand{\Z}{\mathbb{Z}}
\newcommand{\Q}{\mathbb{Q}}
\newcommand{\R}{\mathbb{R}}
\newcommand{\C}{\mathbb{C}}
\newcommand{\F}{\mathbb{F}}
\newcommand{\A}{\mathbb{A}}
\newcommand{\PP}{\mathbb P}
\newcommand{\E}{\mathbb{E}}
\newcommand{\eps}{\varepsilon}
\newcommand{\dd}{\,\mathrm d}
\newcommand{\sref}[2]{\hyperlink{source:#1:#2}{[S#2]}}
\newcommand{\eref}[2]{\hyperlink{extra:#1:#2}{[E#2]}}
% Turn the next switch off for a shorter reading copy. The quotations preserve
% every exactTarget field, including qualifications omitted from a short summary.
\newif\ifshowcatalogue
\showcataloguetrue
\newcommand{\cataloguescope}[1]{\ifshowcatalogue
 \paragraph{Exact scope in the supplied catalogue.}
 \begin{quote}\small #1\end{quote}\fi}

% Research additions dated 27 September 2026; compile twice with LuaLaTeX.
\numberwithin{equation}{section}
\begin{document}
\setcounter{section}{48}
% ================================================================
% problemId: problem.grothendieck-period-conjecture
% source releaseRank: 49; source releaseStatus: open
% exactTarget SHA-256: cd0e5188a2b3b6c66ad1605cc1ac01c150703ecfb43eb476d49db2fc6477edc7
\clearpage
\section{\texorpdfstring{Grothendieck Period Conjecture}{Grothendieck Period Conjecture}}\label{49}
\subsection{Definitions and mathematical statement}
For a Nori motive $M/\Q$, its Betti and de Rham realizations are finite-dimensional rational vector spaces compared over $\C$. Their comparison matrix, in rational bases, supplies the periods of $M$. Its motivic Galois group $G_{\rm mot}(M)$ is the affine algebraic group of tensor automorphisms of the fiber functor on the tensor category generated by $M$. The conjecture is
\[
 \operatorname{trdeg}_{\Q}\Q(\text{entries of the period matrix of }M)
 =\dim G_{\rm mot}(M).
\]
Changes of rational bases do not change the generated period field. The Nori-motive convention and the category generated by $M$ are those of the source.
\par\smallskip\textit{Formulation and concept references:} \sref{49}{1}.
\cataloguescope{Prove that for every Nori motive M over Q, trdeg\_Q Q(periods of M) equals dim G\_mot(M), the dimension of its motivic Galois group.}
\subsection{Short English statement}
Are all algebraic relations among a motive's periods exactly those forced by its motivic symmetries?
% BEGIN RESEARCH ADDITION 49
\subsection{Research attempt}

\paragraph{Current frontier.}
The target is the arithmetic statement for Nori motives over $\Q$, not a
statement about period functions over a complex function field. Bakker and
Tsimerman formulate the former as Conjecture~1.1 and prove a functional
analogue in Theorem~1.1, made precise in Theorem~4.25. Their
Sections~4.2--4.3 fix the Nori category and period-field conventions used
here. The difference between the constant and functional settings is
essential. \sref{49}{1}, \sref{49}{2}
Kreutz--Shen--Vial, Corollary~6.17(i), prove the period-conjecture statement
for K3 surfaces of Picard rank $20$; Kawabe's self-product CM Kummer case
is a subcase, as his revised Remark~3.1 acknowledges. Neither result covers
all Nori motives. Statements for motivated classes or for de Rham--Betti
classes must also be distinguished from the full arithmetic equality.
\eref{49}{1}, \sref{49}{3}
Bertolin's work on $1$-motives explicitly separates general conjectures
of Grothendieck--Andr\'e type from the CM/torsion cases treated there.
\eref{49}{3}

\paragraph{Attack route.}
Three routes deserve separate tests. One can try to make the ideal of
numerical period relations invariant under the motivic group; transitivity
would then force the desired genericity. One can try to transfer functional
transcendence to every arithmetic specialization, but this needs a
specialization theorem with exceptional fibers controlled. Finally, one
can attempt a bounded-degree search for all period relations, which needs
an a priori relation-degree bound, not just Noetherianity. The calculations
below expose the exact gaps in these proposals. They also test the desired
invariance on a small logarithmic motive where the outstanding arithmetic
independence is already substantial.

\paragraph{Attempt.}
\textbf{1. A component-aware genericity criterion.}
Let $G=G_{\rm mot}(M)$ and let
\[
 \mathcal P=\operatorname{Isom}^{\otimes}
       (\omega_{\rm dR},\omega_B)\big|_{\langle M\rangle}
\]
be the comparison torsor. These are the fiber functors on the tensor
category generated by $M$, including duals and subquotients. The comparison
isomorphism determines a complex point $p\in\mathcal P(\C)$.
The torsor formalism and its relation to the period field are those in
\sref{49}{2}, Sections~4.2, 4.3, and 4.6. The following commutative-algebra
argument is supplied here to track the direction of the implication.

Writing $A=\Q[\mathcal P]$, evaluation defines a prime ideal
$I=\ker(A\to\C)$. Coordinates of $\mathcal P$ are generated by matrix
entries of the comparison on $M$ and the inverse determinant, subject to
the tensor relations. Inverting a nonzero determinant does not change the
field generated by its entries. Thus
\[
 \operatorname{Frac}(A/I)=\Q(\mathrm{periods}(M)),\qquad
 \operatorname{trdeg}_{\Q}\Q(\mathrm{periods}(M))
       =\dim(A/I)\le\dim\mathcal P=\dim G.
 \tag{49.1}
\]
This proves the easy upper bound, not the required lower bound.

There is a useful precaution when $G$ is not connected. Choose a finite
number field $K\subset\C$ over which the geometric components of
$\mathcal P$ are defined, and let $C$ be the component containing $p$.
In characteristic zero the torsor is smooth; its geometric connected
components are integral, all of dimension $\dim G$. Hence $K[C]$ is a
domain, and, putting $J=\ker(K[C]\to\C)$, we obtain
\[
 \operatorname{trdeg}_{\Q}\Q(\mathrm{periods}(M))=\dim G
       \quad\Longleftrightarrow\quad J=0.
 \tag{49.2}
\]
Indeed, adjoining $K$ changes no transcendence degree; a nonzero prime
ideal in the coordinate ring of an integral finite-type variety cuts out
a proper subvariety of strictly smaller dimension. The fraction field
of the image here is $K(\mathrm{periods}(M))$. A formulation using
injectivity on the whole coordinate ring without addressing components
would require an extra irreducibility justification.

\textbf{2. The proposed group-invariance step.}
Suppose the $K$-Zariski closure $Z$ of $p$ inside $C$ were stable under the
connected group acting transitively on $C$. After extension to $\C$, choose
any point of the nonempty closed set $Z_{\C}$. Its group orbit is all of
$C_{\C}$, so $Z_{\C}=C_{\C}$, and therefore $J=0$. In ideal language, a
sufficient missing lemma is that the numerical relation ideal $J$ is a
subcomodule for this group action.

No proof of that lemma is obtained. A motivic automorphism acts on the
\emph{space of tensor-compatible comparison isomorphisms}; it has not
thereby been shown to induce an automorphism of the field of their
\emph{actual complex values}. To infer $R(gp)=0$ from $R(p)=0$ would be
to assume the very stability being sought. The torsor gives an algebraic
space in which the periods lie, but it does not make the distinguished
complex point automatically generic. For a transitive component, this
stability condition is equivalent to genericity itself, not a known
independently weaker arithmetic theorem.

\textbf{3. A finite-degree obstruction to a shortcut.}
The following exact model shows why checking a large finite collection
of tensor or polynomial relations cannot be promoted to unrestricted
algebraic independence without a further theorem. Fix $D\ge1$, choose
any transcendental complex number $\tau$, and set
\[
 T=\mathbb G_m^2,\qquad
 p_D=(\tau,\tau^{2D+1}).
\]
For Laurent monomials $X^aY^b$ with $|a|,|b|\le D$, evaluation gives
$\tau^{a+(2D+1)b}$. These exponents are all distinct: equality would
imply
\[
 (2D+1)(b-b')=a'-a,
\]
whose right side has absolute value at most $2D$, forcing $b=b'$ and then
$a=a'$. Consequently no nonzero Laurent polynomial supported in that box
vanishes at $p_D$, even with coefficients in a fixed number field. Multiply
by a sufficiently large power of $\tau$ and use its transcendence to see
this explicitly. Nevertheless
\[
 Y-X^{2D+1}=0\quad\hbox{at }p_D,\qquad
 \operatorname{trdeg}_{\Q}\Q(p_D)=1<2=\dim T.
 \tag{49.3}
\]
The model passes every such degree-$D$ test and still fails genericity.
It is \emph{not} asserted to be the comparison torsor and period point
of any Nori motive. It falsifies only a general algebraic inference used
in the proposed proof method. Noetherianity says that an already given
ideal has finitely many generators; it supplies neither their numerical
identification nor the needed uniform degree bound.

\textbf{4. What functional transcendence does and does not supply.}
For a basic test, the germ $\log t$ is transcendental over $\C(t)$.
If it were algebraic, continuation around a base point outside the finite
branch locus of its polynomial equation could produce only finitely many
values. Continuation around $0$ instead gives $\log t+2\pi i k$ for all
integers $k$, a contradiction. Yet the branch with $\log1=0$ has an
algebraic value at $t=1$.

This is an obstruction to the inference ``a function is transcendental,
therefore every specified algebraic specialization has the full expected
transcendence degree.'' It is not a counterexample to the arithmetic period
conjecture. In particular, the simple relative-pair representation using
marked points $\{1,t\}$ changes when the points collide at $t=1$; one may
not silently preserve that motive's group dimension at the collision.
The example isolates the missing specialization control, rather than
claiming that a theorem about nondegenerate families has been contradicted.
Bakker--Tsimerman's functional theorem is used with its stated quantifiers,
not as such an arithmetic specialization theorem. \sref{49}{2}

\textbf{5. An actual logarithmic motive as an arithmetic stress test.}
Take the Nori motive
\[
 M_0=\mathcal H^1(\mathbb G_m,\{1,2,3\}).
\]
Its period field is $\Q(2\pi i,\log2,\log3)$, and its motivic Galois
group is an extension of $\mathbb G_m$ by $\mathbb G_a^2$, hence has
dimension $3$. This is the case $\alpha_1=2,\alpha_2=3$ of the explicit
calculation in \sref{49}{2}, Example~4.4. Thus this particular instance
requires
\[
 \operatorname{trdeg}_{\Q}\Q(2\pi i,\log2,\log3)=3.
 \tag{49.4}
\]
The elementary linear check does go through. If
$a(2\pi i)+b\log2+c\log3=0$ with rational coefficients, imaginary parts
give $a=0$. Clearing denominators and exponentiating gives $2^B3^C=1$;
unique prime factorization forces $B=C=0$. This proves linear, not
algebraic, independence.

Schanuel's conjecture applied to the three linearly independent numbers
$2\pi i,\log2,\log3$ would imply (49.4), since their exponentials are
$1,2,3$. This flags a dependency on the separate unresolved conjecture,
not a proof of (49.4). Conversely, nothing here claims that the original
$\Q$-motive period conjecture implies full Schanuel for arbitrary complex
numbers. The broader Grothendieck--Andr\'e conjecture over arbitrary
subfields of $\C$ is a different statement; the distinction is explicit
in \sref{49}{2}, Section~4.3, and \eref{49}{3}, Section~6.

\textbf{6. A sector where both sides can actually be calculated.}
For comparison, let $M$ be a \emph{direct sum} of an Artin motive and
finitely many pure Tate motives $\Q(m_j)$. This does not include arbitrary
extensions of these objects. Artin periods are algebraic, while the Tate
periods are integral powers of $2\pi i$ or their inverses. If every
$m_j=0$, the period field is algebraic and the motivic group is finite,
so both quantities are zero. If some $m_j\ne0$, the group has a
one-dimensional torus image, with only a finite Artin part, and the period
field has transcendence degree exactly one. Here the lower bound follows
from the transcendence of $\pi$. \eref{49}{2}
These are familiar special cases, included as a check of dimensions,
finite components, and twists, not as a newly discovered theorem.

\paragraph{Stress test.}
The upper bound (49.1) survives changes of rational bases and adjoining a
finite coefficient field. The component argument does not assume a rational
point on the torsor. The proposed invariance lemma, however, is not supplied
by the definition of a motivic automorphism. The finite-degree model (49.3)
shows exactly how hidden nonlinear relations can escape every prescribed
finite test; it says nothing against the additional geometric structure
of genuine motives. Likewise, the logarithm specialization is a test of
an inference, not a motive with a verified dimension discrepancy.

No numerical independence computation is reported as evidence for a
proof: a small nonzero residual cannot establish algebraic independence,
and an approximate integer relation cannot establish an exact one. The
proof of linear independence in the genuine test (49.4) does not control
quadratic or higher relations. Extending the Artin--Tate calculation to
nontrivial mixed extensions already encounters such logarithmic periods.

\paragraph{Outcome.}
\textbf{Outcome: Exploratory approach with unresolved gap.}

The attempt supplies a component-aware genericity formulation, spells out
the exact unproved group-invariance step, and proves a finite-degree
countermodel to a possible shortcut. It also reduces a concrete logarithmic
motive to the explicit algebraic-independence requirement (49.4) and checks
a known direct-sum Artin--Tate sector. It produces neither a general
arithmetic lower bound nor a Nori-motive counterexample. A resolution by
this route needs a genuinely arithmetic theorem preventing proper
numerical relation loci, not just the formal symmetries of the torsor,
functional transcendence, or a finite number of numerical tests. No claim
of historical novelty is made for these auxiliary deductions.
% END RESEARCH ADDITION 49
\subsection{Sources}
\begin{itemize}\raggedright
\item[\textnormal{[S1]}]\hypertarget{source:49:1}{} Formal statement, Conjecture 1.1.
\url{https://www.cambridge.org/core/journals/forum-of-mathematics-sigma/article/functional-transcendence-of-periods-and-the-geometric-andregrothendieck-period-conjecture/6B3E39BB9BC68DC2BC6E793D15B0735B}.
{\small\textit{Catalogue formulation source.}}
\item[\textnormal{[S2]}]\hypertarget{source:49:2}{} Functional Transcendence of Periods and the Geometric André--Grothendieck Period Conjecture.
\url{https://arxiv.org/html/2208.05182v2}.
{\small\textit{Catalogue research/status reference; not a proof endorsement.}}
\item[\textnormal{[S3]}]\hypertarget{source:49:3}{} Grothendieck's period conjecture for Kummer surfaces of self-product CM type.
\url{https://arxiv.org/abs/2303.05030}.
{\small\textit{Catalogue research/status reference; not a proof endorsement.}}
% BEGIN RESEARCH SOURCES 49
\item[\textnormal{[E1]}]\hypertarget{extra:49:1}{} Tobias Kreutz, Mingmin Shen, and Charles Vial, \emph{De Rham--Betti classes with coefficients}, arXiv:2206.08618v3, especially Corollary 6.17(i) and Remark 6.18.
\url{https://arxiv.org/html/2206.08618v3}.
{\small\textit{Added primary source; the Picard-rank-20 K3 result is distinguished from the general conjecture. Consulted 27 September 2026.}}
\item[\textnormal{[E2]}]\hypertarget{extra:49:2}{} Ferdinand Lindemann, \emph{Ueber die Zahl $\pi$}, Mathematische Annalen 20 (1882), 213--225, DOI:10.1007/BF01446522.
\url{https://doi.org/10.1007/BF01446522}.
{\small\textit{Classical transcendence of $\pi$, used only for the Tate-sector check. Publisher bibliographic record checked 27 September 2026; the subscription full text was not inspected.}}
\item[\textnormal{[E3]}]\hypertarget{extra:49:3}{} Cristiana Bertolin, \emph{Schanuel conjecture for 1-motives}, arXiv:2509.08700v1, Sections 6--7.
\url{https://arxiv.org/html/2509.08700v1}.
{\small\textit{Added primary source; the broader Grothendieck--Andr\'e formulation and the restricted CM/torsion result, not a proof of the present universal statement. Consulted 27 September 2026.}}
% END RESEARCH SOURCES 49
\end{itemize}

\end{document}
